\origpage{146}{ROBERTSON AND SEYMOUR}
\hspace*{12pt}(i) \textit{$\mathcal{T}^*$ controls a $K_p$ minor of $G$, or}

\hspace*{12pt}(ii) \textit{there is a $\Sigma'$-span of order $\ge \phi$, where $\Sigma'$ is a surface obtained by adding a crosscap to $\Sigma$, or}

\hspace*{12pt}(iii) \textit{no $\Sigma$-span of order $\ge \theta$ is $(\lambda, \mu)$-flat.}

\textit{Proof.} If the result holds for $4, \phi, \mu$ then it holds for $p, \phi, \mu$ for all $p \le 4$, because if $\mathcal{T}^*$ controls a $K_4$ minor then it controls a $K_p$ minor for all $p \le 4$. Hence we may assume that $p \ge 4$. Let $\kappa, \phi_1$ satisfy (8.1) (with $\phi$ replaced by $\phi_1$), and let $\phi_2 = \max(\phi, \phi_1)$. Let $\theta, \lambda, \rho$ satisfy (1.1) (with $\phi$ replaced by $\phi_2$). We may assume that $\theta \ge p$. Then we claim $\kappa, \theta, \lambda, \rho$ satisfy (8.2). For let $\mathcal{T}^*, G$ be as in the theorem. If $\mathcal{T}^*$ has order $< p$ then (8.2)(iii) holds, since $\theta \ge p$. Thus we may assume that $\mathcal{T}^*$ has order $\ge p$. By (1.1), one of (1.1)(i), (ii), (iii) holds. If (1.1)(i) holds then (8.2)(i) holds, by (8.1), and if (1.1)(ii) or (iii) holds, then (8.2)(ii) or (iii) holds. $\square$

One of the conclusions of (8.2) is that no $\Sigma$-span of order $\ge \theta$ is $(\lambda, \mu)$-flat, and the second modification we wish to make concerns this. If $\mathcal{T}^*$ is a tangle in $G$ and $H, \eta, \mathcal{T}$ is a $\Sigma$-span of order $\theta$, where $\Sigma$ is a surface with $\text{bd}(\Sigma) = \emptyset$, we say that $H, \eta, \mathcal{T}$ is $(\lambda, \mu)$-level if

\hspace*{12pt}(i) $\theta \ge 4\lambda + \mu + 2$

\hspace*{12pt}(ii) $H, \eta, \mathcal{T}$ is not $\mu$-stepped, and

\hspace*{12pt}(iii) for every $\Sigma$-span $H', \eta', \mathcal{T}'$, obtained from $H, \eta, \mathcal{T}$ by rearranging within $\lambda$ of some $z \in A(H)$, there is no $\eta'(H')$-path in $G$ with ends $\eta'(u), \eta'(v)$ with $d'(u, v) \ge \mu$ and $d(z, u), d(z, v) \le \lambda$.

(The difference between (i)--(iii) here and the definition of $(\lambda, \mu)$-flat is the requirement that $d(z, u), d(z, v) \le \lambda$.)

\stmt
(8.3) \textit{Let $\lambda, \mu \ge 0$ be integers, let $\mathcal{T}^*$ be a tangle in a graph $G$, and let $\Sigma$ be a surface with $\text{bd}(\Sigma) = \emptyset$. Then every $\Sigma$-span which is $(\lambda + \mu + 2, \mu)$-level is $(\lambda, \mu)$-flat.}

\textit{Proof.} Suppose that $H, \eta, \mathcal{T}$ is a $\Sigma$-span of order $\theta$ is not $(\lambda, \mu)$-flat; we shall show that it is not $(\lambda + \mu + 2, \mu)$-level. If $\theta < 4(\lambda + \mu + 2) + \mu + 2$ then $H, \eta, \mathcal{T}$ is not $(\lambda + \mu + 2, \mu)$-level as required, and so we assume that $\theta \ge 4\lambda + 5\mu + 10$. Since $H, \eta, \mathcal{T}$ is not $(\lambda, \mu)$-flat and $\theta \ge 4\lambda + \mu + 2$, there is a $\mu$-stepped $\Sigma$-span $H', \eta', \mathcal{T}'$ of order $\theta - 4\lambda - 2$, obtained from $H, \eta, \mathcal{T}$ by rearranging within $\lambda$ of some $z \in A(H)$. Let $P$ be an $\eta'(H')$-path with ends $\eta'(a), \eta'(b)$, such that $d'(a, b) \ge \mu$.
